\documentclass[11pt]{article}
\usepackage[utf8]{inputenc}
\usepackage[T1]{fontenc}
\usepackage{amsmath,amssymb}
\usepackage{graphicx}
\usepackage{booktabs}
\usepackage{hyperref}
\usepackage[margin=1in]{geometry}
\usepackage{natbib}
\bibliographystyle{unsrtnat}

\title{Time-, Tissue- and Treatment-Associated Heterogeneity in Tumour-Residing Migratory DCs}
\author{Authors}
\date{}

\begin{document}
\maketitle

\begin{abstract}
Tumour dendritic cells (DCs) internalise antigen and upregulate CCR7, which directs their migration to tumour-draining lymph nodes (dLN). CCR7 expression is coupled to a maturation programme enriched in regulatory molecule expression, including PD-L1, termed mRegDC. However, the spatio-temporal dynamics and role of mRegDCs in anti-tumour immune responses remain unclear. Using photoconvertible mice to precisely track DC migration, we found that mRegDCs were the dominant DC population arriving in the dLN, but a subset remained tumour-resident despite CCR7 expression. These tumour-retained mRegDCs were phenotypically and transcriptionally distinct from their dLN counterparts and were heterogeneous. Specifically, they demonstrated a progressive reduction in the expression of antigen presentation and pro-inflammatory transcripts with more prolonged tumour dwell-time. Tumour mRegDCs spatially co-localised with PD-1$^+$CD8$^+$ T cells in human and murine solid tumours. Following anti-PD-L1 treatment, tumour-residing mRegDCs adopted a state enriched in lymphocyte stimulatory molecules, including OX40L, which was capable of augmenting anti-tumour cytolytic activity. Altogether, these data uncover previously unappreciated heterogeneity in mRegDCs that may underpin a variable capacity to support intratumoural cytotoxic T cells, and provide insights into their role in cancer immunotherapy.
\end{abstract}

\section{Introduction}

Dendritic cells (DCs) capture tumour antigens and upregulate the chemokine receptor CCR7, which directs their migration to secondary lymphoid organs where they present antigens to T cells \citep{Wculek2019,Roberts2016,Sallusto1998}. Two major conventional DC (cDC) subsets have been identified in tumours: cDC1, which specialise in cross-presenting tumour antigens to CD8$^+$ T cells and are associated with improved survival \citep{Bottcher2018}, and cDC2, which are involved in CD4$^+$ T cell activation.

The recent application of single-cell technologies to tissue samples has enabled the distinction of an intratumoral DC population that may arise from both cDC1 and cDC2 but demonstrates a conserved programme of maturation-associated gene expression. This DC subset or state has been variously labelled as ``migratory DCs'', ``mature DCs enriched in immunoregulatory molecules'' (mRegDCs), and ``LAMP3$^+$ DCs'' \citep{Maier2020,Mulder2021,Zilionis2019}. Regardless of nomenclature, single-cell atlasing efforts have identified mRegDCs in multiple human cancers \citep{Cheng2021,Zhang2019,Maier2020}, and the acquisition of this maturation programme appears dependent on the uptake of tumour antigens \citep{Maier2020}.

The presence of mature LAMP$^+$ DCs within lymphoid follicles has been associated with improved prognosis in non-small cell lung cancer \citep{DeCastro2012,Dieu-Nosjean2008,Goc2014}. However, the functional significance of mRegDCs in anti-tumour immunity remains debated. The mRegDC programme includes co-expression of both immunostimulatory (e.g., CD80, CD86) and immunoregulatory (e.g., PD-L1, PD-L2) molecules, raising the question of whether mRegDCs promote or suppress anti-tumour T cell responses. Furthermore, PD-L1 expression by DCs has been shown to be essential for effective responses to anti-PD-L1 antibodies \citep{Oh2020,Peng2020}.

Here, we use photoconvertible Kaede transgenic mice to precisely track the spatio-temporal dynamics of mRegDCs in tumours and draining lymph nodes. We reveal that mRegDCs are transcriptionally heterogeneous, with tumour-retained mRegDCs being phenotypically distinct from those that successfully migrate to the dLN. Prolonged tumour residence is associated with progressive loss of antigen presentation capacity. Critically, anti-PD-L1 treatment enriches for an activated mRegDC state characterised by expression of lymphocyte co-stimulatory molecules including OX40L, which functionally augments anti-tumour CD8$^+$ T cell responses.

\section{Results}

\subsection{mRegDC Signatures Are Associated with Improved Survival in Human Cancers}

We analysed the transcriptomes of 4,045 human solid tumours from the Cancer Genome Atlas (TCGA) \citep{TCGA2013}. mRegDC gene signature enrichment was associated with improved overall survival across multiple cancer types, including melanoma, colorectal adenocarcinoma, breast cancer, and lung adenocarcinoma (Extended Data Figure 1A). Further analysis of 1,853 human breast tumours from METABRIC revealed cancer subtype-specific associations with survival. These data suggest that mRegDCs contribute to anti-tumour responses across a range of human cancers.

\subsection{mRegDCs Are Transcriptionally Heterogeneous and Some Remain Within the Tumour}

To investigate the mechanisms by which mRegDCs promote anti-tumour immunity, we characterised their spatio-temporal dynamics using photoconvertible Kaede transgenic mice bearing MC38, MC38-Ova, and CT26 subcutaneous tumours. Tumours were transcutaneously photoconverted on day 13, converting all infiltrated cells from default Kaede-green to Kaede-red fluorescence \citep{Tomura2008}. After 24--72h, tumours were harvested, enabling distinction of newly infiltrating Kaede-green cells from tumour-retained Kaede-red cells.

Single-cell RNA sequencing of sorted CD45$^+$ cells generated 80,556 transcriptomes including 32,191 myeloid cells. Unbiased clustering of DCs revealed 8 distinct clusters, including cycling DCs, cDC1, cDC2s, and mRegDCs (Figure~\ref{fig:landscape}B--D). mRegDC showed high expression of \textit{Ccr7}, \textit{Cd274}, and \textit{Pdcd1lg2}, and expressed higher levels of surface PD-L1 than other immune cells.

Surprisingly, tumour mRegDCs were predominantly Kaede-red, indicating they had resided in the tumour for over 48h (Figure~\ref{fig:landscape}F). In contrast, cDC1 and cDC2 were mostly Kaede-green, consistent with newly infiltrating cDCs differentiating into mRegDCs following uptake of tumour antigen. Immunofluorescence microscopy confirmed the presence of Kaede-red MHC-II$^+$CCR7$^+$ cells within the tumour.

Pseudotime analysis revealed a trajectory from cDC2, through an intermediate mRegDC\_1 state, to terminal mRegDC\_2 and mRegDC\_3 states (Figure~\ref{fig:landscape}I--J). RNA velocity analysis confirmed this cell-state transition. The receptor tyrosine kinase Axl, which recognises apoptotic cells and induces the mRegDC programme, was upregulated along this trajectory (Figure~\ref{fig:landscape}K).

\subsection{Migrated mRegDCs Are Phenotypically Distinct from Tumour-Residing Populations}

Kaede-red CCR7$^+$PD-L2$^+$ DCs were readily detectable in the dLN but not in the contralateral non-draining LN (Figure~\ref{fig:migration}A). Among Kaede-red DCs in the dLN, essentially all were mRegDC (Figure~\ref{fig:migration}B--C). scRNA-seq confirmed that 94\% of Kaede-red myeloid cells in the dLN were mRegDCs, establishing mRegDCs as the dominant myeloid cell tumour emigrants arriving in the dLN.

GSEA revealed that tumour mRegDCs and dLN mRegDCs were enriched for distinct gene programmes. CD80 and CD86 were higher in tumour-residing mRegDCs, while ICOSL was significantly higher in migrated mRegDCs (Figure~\ref{fig:migration}G--H). \textit{Il12b}, which drives anti-tumour cytotoxic T cell activity, was enriched in tumour mRegDCs, while \textit{Il15ra} was higher in mRegDC in the dLN.

Strikingly, over 80\% of Kaede-red DCs from the dLN mapped to the tumour mRegDC\_1 cluster (Figure~\ref{fig:migration}F). Only 9\% resembled mRegDC\_2/3, despite these being the dominant tumour-resident mRegDC states. These data suggest that mRegDC\_1, which have most recently acquired the mRegDC programme, are the main population to seed the dLN, while mRegDC\_2 and mRegDC\_3 are terminal, tumour-residing states.

\subsection{Tumour-Retained mRegDCs Progress Towards an ``Exhausted''-Like State}

During the progression from mRegDC\_1 to mRegDC\_3, there was a decrease in MHC-II expression, including the class-II invariant chain \textit{Cd74} (Figure~\ref{fig:exhaustion}A). DC antigen presentation genes decreased as cells transitioned towards tumour-retained states (Figure~\ref{fig:exhaustion}B--C), including molecules involved in antigen processing (\textit{Ctsb}, \textit{Ctsl}, \textit{Ifi30}), transport (\textit{Tap1}, \textit{Tap2}), and MHC loading (\textit{H2-DMa}, \textit{H2-Oa}, \textit{Tapbp}, \textit{Calr}).

Genes involved in innate immune function and response to inflammatory stimuli, including ``interferon gamma/alpha response'', ``inflammatory response'', and ``Toll-like receptor signalling'' pathways, significantly decreased as cells transitioned from mRegDC\_1 to mRegDC\_3 (Figure~\ref{fig:exhaustion}D--E). Specifically, mRegDC\_3 showed reduced expression of innate immune response genes (\textit{Nlrp3}, \textit{Tlr1}, \textit{Cd14}), DC activation and migration genes (\textit{Axl}, \textit{Ccr7}, \textit{Rgs1}), lymphocyte co-stimulation (\textit{Cd40}, \textit{Tnfsf9}, \textit{Pvr}), and cytokines (\textit{Cxcl9}, \textit{Il1b}).

\subsection{Anti-PD-L1 Treatment Enriches for an Activated mRegDC State}

Assessment of differential abundance showed significant enrichment within mRegDC\_2 neighbourhoods following anti-PD-L1 treatment, but relative depletion of mRegDC\_3 (Figure~\ref{fig:exhaustion}F). The transcriptome of mRegDC\_2 was highly activated compared to mRegDC\_3, with increased expression of immune signalling genes including \textit{Cd40}, \textit{Il1b}, \textit{Tnf}, and \textit{Nfkbia} (Figure~\ref{fig:exhaustion}G).

We identified several lymphocyte activation ligands significantly upregulated on mRegDC\_2 versus other tumour mRegDC states and dLN mRegDCs, including \textit{Il12b}, \textit{Tnfsf4} (OX40L), \textit{Tnfsf9} (4-1BBL), \textit{Cd70}, and \textit{Pvr}. Using OX40L reporter mice, we confirmed that tumour mRegDCs express OX40L, but it was not expressed by tumour cDCs or mRegDCs in the dLN (Figure~\ref{fig:exhaustion}H).

\subsection{mRegDCs Interact with CD8$^+$ T Cells in Human and Murine Tumours}

Deconvolution of 521 human CRC transcriptomes revealed the highest correlation between mRegDC and CD8$^+$ T cell transcripts (Figure~\ref{fig:interaction}A), replicated in melanoma, breast, and lung tumour biopsies. CellPhoneDB analysis showed that mRegDCs had stronger predicted interactions with effector CD8$^+$ T cells than other myeloid cells, including via CD28, CTLA-4, and PD-1 engagement (Figure~\ref{fig:interaction}C).

Spatial transcriptomics of human CRC, melanoma, and breast tumours revealed hotspots with co-localised expression of mRegDC and effector CD8$^+$ T cell genes (Figure~\ref{fig:interaction}D--F). In PIC-seq data from NSCLC, mRegDC-CD8$^+$ T cell doublets were more frequent than other mRegDC-T cell combinations, and these PICs highly co-expressed \textit{CD8B}, \textit{PRF1}, \textit{GZMB}, and \textit{PDCD1} (Figure~\ref{fig:interaction}G).

\subsection{Anti-PD-L1 Promotes Immunogenic mRegDC--CD8$^+$ T Cell Interactions}

Among CD8$^+$ T cells in tumours, we identified a major Prf1$^{high}$Gzmb$^{high}$Pdcd1$^+$Havcr2$^+$ cluster resembling ``exhausted'' T (T\textsubscript{EX}) cells and a Pdcd1$^+$Tcf7$^+$Slamf6$^+$ cluster resembling ``stem-like'' T cells (Figure~\ref{fig:treatment}A--B). CD8$^+$ T\textsubscript{EX} cells showed the largest transcriptional response to anti-PD-L1, including upregulation of TCR signalling and IL2-STAT5 signalling genes (Figure~\ref{fig:treatment}C).

Cell-cell communication analysis identified TNF-superfamily interactions -- TNFSF4--TNFRSF4 (OX40L--OX40), TNFSF9--TNFRSF9 (4-1BBL--4-1BB), CD70--CD27, and PVR-mediated interactions -- upregulated between mRegDC\_2 and CD8$^+$ T\textsubscript{EX} cells (Figure~\ref{fig:treatment}D). Immunofluorescence confirmed that mRegDCs and proliferating 4-1BB$^+$ CD8$^+$ T cells co-localised within tumours (Figure~\ref{fig:treatment}E).

Functional co-culture of tumour-antigen-experienced mRegDC-like BMDCs with OT-I CD8$^+$ T cells showed that anti-PD-L1 antibodies increased activated, clonally expanded OT-I cells and enhanced granzyme B production (Figure~\ref{fig:treatment}H--I). OX40L-deficient mReg-BMDCs showed reduced capacity to activate OT-I cells (Figure~\ref{fig:treatment}J). Recombinant IFN$\gamma$ treatment of mReg-BMDCs upregulated OX40L, PVR, and CD40, increasing T cell activation in an OX40L-dependent manner (Figure~\ref{fig:treatment}K).

\subsection{Conserved mRegDC Heterogeneity and CD8$^+$ T Cell Crosstalk in Human Cancers}

In scRNA-seq of human CRC, mRegDCs showed a gradient of MHC-II expression (Figure~\ref{fig:human}A--B). Markers upregulated on tumour-residing MHC-II$^{low}$ mRegDCs in murine tumours were also preferentially expressed in human MHC-II$^{low}$ mRegDCs, including IL15, PVR, TNFSF4, TNFSF9, and CD70.

Integration of bulk transcriptomes from the IMvigor210 trial for metastatic urothelial carcinoma with scRNA-seq revealed that myeloid cells associated with favourable clinical response expressed transcripts associated with tumour-residing mRegDCs, including \textit{CCR7}, \textit{CD274}, \textit{IL15}, \textit{PVR}, and \textit{CD70} (Figure~\ref{fig:human}C--E). In breast cancer, mRegDCs from tumours with T cell clonotype expansion following anti-PD-1 treatment were enriched in transcripts associated with the ICB-induced mRegDC\_2 cluster (Figure~\ref{fig:human}F).

\section{Discussion}

Our study reveals previously unappreciated heterogeneity in tumour-residing mRegDCs that has implications for understanding anti-tumour immunity and responses to immune checkpoint blockade. Using photoconvertible mice to track DC dynamics in real time, we demonstrate that mRegDCs are the dominant DC population that seeds the dLN but that a substantial proportion remain resident in the tumour despite CCR7 expression.

Tumour-retained mRegDCs undergo a progressive transition characterised by reduced expression of antigen presentation machinery, pro-inflammatory genes, and co-stimulatory molecules, resembling the ``exhaustion'' phenotype described in T cells. This time-associated heterogeneity may explain conflicting reports about whether mRegDCs are pro- or anti-tumoural: recently formed mRegDCs (mRegDC\_1) retain immunostimulatory capacity and seed the dLN, while terminally differentiated tumour-resident states (mRegDC\_2/3) show reduced effector function but maintain high PD-L1/L2 expression.

Critically, anti-PD-L1 treatment redirects the mRegDC trajectory towards the activated mRegDC\_2 state, enriched in lymphocyte co-stimulatory molecules including OX40L, 4-1BBL, CD70, and PVR. The OX40L--OX40 axis in particular appears functionally important, as OX40L-deficient DCs showed reduced capacity to activate antigen-specific CD8$^+$ T cells. These findings provide mechanistic insight into how anti-PD-L1 therapy modulates DC function within the tumour microenvironment and suggest that the mRegDC--CD8$^+$ T cell axis is a key determinant of immunotherapy efficacy.

The conservation of mRegDC heterogeneity and CD8$^+$ T cell interactions across human cancers, including CRC, melanoma, breast cancer, lung cancer, and urothelial carcinoma, underscores the translational relevance of these findings and identifies potential biomarkers and therapeutic targets for improving cancer immunotherapy.

\section{Methods}

\subsection{Mouse Models and Tumour Experiments}

C57BL/6 Kaede, BALB/c Kaede, OX40L reporter, OX40L$^{fl/fl}$, and CD11c$^{cre}$ OX40L$^{fl/fl}$ mice were maintained at the University of Birmingham. MC38, MC38-Ova, and CT26 cells ($2.5 \times 10^5$) were injected subcutaneously. Tumours were photoconverted transcutaneously on day 13. Anti-PD-L1 antibodies were administered on days 7, 10, and 13.

\subsection{Single-Cell RNA Sequencing}

CD45$^+$TER119$^-$ cells were sorted by FACS into CD11b$^+$ and CD11b$^{-/low}$ fractions, combined for scRNA-seq analysis, and annotated based on transcriptomes. Quality control yielded 80,556 single cell transcriptomes. Pseudotime, RNA velocity, PAGA, and differential abundance analyses were performed using standard pipelines.

\subsection{Functional Assays}

BMDCs were cultured with UV-irradiated MC38-Ova cells to generate mRegDC-like cells. OT-I CD8$^+$ T cells were co-cultured with mReg-BMDCs for 3 days with or without anti-PD-L1 antibodies. OX40L-deficient mReg-BMDCs were generated from CD11c$^{cre}$ Tnfsf4$^{fl/fl}$ mice.

\bibliography{references}

\end{document}
